\subsection{A uniform mass bound for active prefixes}
\label{sec:saturation-mass-bound}
Let $S_w$ be the union of the expanded intervals at width $w$, and write
$Q_w=|S_w|$. List these active positions in decreasing score order, breaking
numerical ties arbitrarily but consistently. Numerical equality here is
that of the validated binary64 scores, including ties created by input
conversion. Only ends of complete active-score ties are valid event indices;
write their positive prefix lengths as
\[
 0=b_0<b_1<\cdots<b_g=Q_w.
\]
For $1\le q\le Q_w$, define the tie-rounding map
\[
 \rho_w(q)=\min\{b_j:b_j\ge q\}.
\]
Global positions outside $S_w$ may share these scores. They remain included
in the event's global prediction count, but are not counted in $Q_w$ or $b_j$.

\begin{lemma}[Active support and prefix mass]
For the stated half-width square-root ramps, $S_w$ is precisely the set of
positions with positive label weight and precisely the set of positions
capable of hitting an expanded component. Every active original position
contributes one to $A+B/2$; every active non-original position contributes
at least $1/(2\sqrt2)$. Therefore, after any $q$ active positions have been
admitted, even within a score tie,
\[
 A+B/2\ge q/(2\sqrt2).
\]
\end{lemma}
\begin{proof}
At widths zero and one the active positions are exactly the originals, each
of weight one. At width $w\ge2$, a non-original position is in $S_w$ exactly
when at least one original interval is at integer distance
$d\le\lfloor w/2\rfloor$. Its contribution from that interval is
$\sqrt{1-d/w}\ge1/\sqrt2$. Summing additional nonnegative contributions
and clipping at one preserves this lower bound. Such a position contributes
half its weight to $A+B/2$, while an original contributes one. Summing gives
the inequality. The same union defines expanded-component membership, so
positions outside $S_w$ cannot change either mass or component hits.
\end{proof}

Let $M_w$ be the first complete active-event prefix length with
$A+B/2\ge P_0$, and let $E_w$ be the first complete active-event prefix
length that hits all $K_w$ expanded components. Both exist: at $b_g=Q_w$,
all originals have been admitted, so $A=P_0$, and every expanded component
contains an admitted original. Put
\[
 m=\lceil2\sqrt2 P_0\rceil.
\]

\begin{proposition}[Mass and coverage bound with complete ties]
For the evaluator stopped at the first mass-and-coverage saturation event,
\begin{align}
 J_w&=\max\{M_w,E_w\},\\
 M_w&\le\rho_w(\min\{Q_w,m\}),\\
 J_w&\le\rho_w\!\left(\min\{Q_w,\max\{E_w,m\}\}\right).
 \label{eq:mass-coverage-bound}
\end{align}
If $\tau_w=\max_j(b_j-b_{j-1})$ is the largest active-score tie size, then
\[
 J_w\le\min\{Q_w,\max\{E_w,m+\tau_w-1\}\}.
\]
\end{proposition}
\begin{proof}
Mass and coverage are each nondecreasing under threshold relaxation.
Their first complete-event indices are $M_w$ and $E_w$, so both conditions
first hold at their maximum. If $Q_w\ge m$, the preceding lemma guarantees
the mass condition after $m$ positions; it still holds after completing
that tie. If $Q_w<m$, the final event contains every original and satisfies
the mass condition. This proves the bound on $M_w$. Since $E_w$ is already
a complete-event prefix length, taking its maximum with the rounded mass
bound gives~\eqref{eq:mass-coverage-bound}. Finally,
$\rho_w(q)\le\min\{Q_w,q+\tau_w-1\}$, which yields the stated additive-tie
bound. No partial tie is used as an evaluation threshold or stopping point.
\end{proof}

\begin{corollary}[One original interval]
If $K=1$ and the validated scores within $S_L$ are pairwise distinct, then
for every score ordering and every width,
\[
 J_w\le\min\{Q_w,\lceil2\sqrt2 P_0\rceil\}.
\]
Under the preceding computation model, the complete scalar evaluator takes
$O(n+(L+1)P_0)=O(n+LP_0)$ time and $O(n+Q)$ auxiliary words. In particular,
a single original anomalous position ($P_0=K=1$) gives $J_w\le3$ and total
$O(n+L)$ time, without a score-quality assumption. At widths zero and one,
$J_w=1$ in that singleton case.
\end{corollary}
\begin{proof}
There is one expanded component at every width. The first active position
hits it, and with distinct active scores $E_w=1$ and $\rho_w(q)=q$.
Apply the proposition. Union-find operations are constant time when $K=1$,
and sum the per-width bound. Since $P_0\le n$, the additional $P_0$ term is
absorbed into $O(n+LP_0)$. For $P_0=1$, $\lceil2\sqrt2\rceil=3$.
\end{proof}

The distinctness requirement concerns active/support scores, not arbitrary
off-support scores. Equal scores exclusively outside $S_L$, or off-support
positions tied with a single active position, do not enlarge an active tie.
More generally, with one original interval and $\tau_w\le\tau$ at every
width, the same analysis gives $O(n+(L+1)(P_0+\tau))$ time.

Arbitrary active ties prevent the singleton corollary for this
position-by-position evaluator. For $L\ge1$ and $n=2L+1$, place the only original at
zero-based position $L$ and assign all $n$ positions the same score. At
width $w$, the sole active event contains
$Q_w=2\lfloor w/2\rfloor+1$ positions. The evaluator visits that complete
event, so $J_w=Q_w$ and $\sum_{w=0}^L J_w=\Theta(L^2)$, although
$P_0=K=1$. This is a counterexample to dropping the tie restriction from
this evaluator's bound; it is not a lower bound for all possible VUS
algorithms or for implementations that aggregate tied masses differently.

For multiple original intervals, the mass bound alone does not bound the
coverage prefix $E_w$. A component whose active scores are all low can keep
$E_w$ near $Q_w$ even after the mass condition holds. Thus neither small
$P_0$ alone nor the mass inequality supplies a general multi-component
$O(n+LP_0)$ guarantee. The proof concerns the pinned ramp, denominator,
component-hit definition, and complete-threshold scalar target in exact
metric arithmetic; finite-precision stopping is numerically validated
separately.

\subsection{Linear input-inspection optimality for a singleton anomaly}
\label{sec:singleton-optimality}
The preceding upper bound admits a matching input-inspection lower bound
on a restricted class. The input-access model matters: an algorithm is
deterministic and adaptively reads individual score coordinates in random
access. Reading a coordinate reveals its validated binary64 value. The labels,
$n$, and $L$ may be supplied without charge, and validity and pairwise score
distinctness are promised. No precomputed ordering, score summary, or
certificate is supplied. A unit-time word operation reads at most a constant
number of score coordinates. Thus the argument does not charge input
validation and does not concern parallel or aggregate-query oracles.

\begin{theorem}[Singleton score-query lower bound]
\label{thm:singleton-query-lower}
Fix integers $n\ge4$ and $0\le L\le n$, and assume that the score domain
contains $n-1$ distinct finite binary64 values below one. Every deterministic
algorithm that exactly evaluates scalar VUS-PR on all inputs with these
$n,L$, one original anomalous position, and pairwise distinct validated
scores must, on some such input, inspect at least
\[
 n-\lfloor L/2\rfloor-1
 \;\ge\;\lceil n/2\rceil-1
 \;\ge\;n/4
\]
distinct score coordinates. The same lower bound applies when the requested
output is the pair VUS-ROC and VUS-PR.
\end{theorem}
\begin{proof}
Fix the label vector $y=(1,0,\ldots,0)$; the algorithm may know this vector
for free. Set $r=\lfloor L/2\rfloor$. Because the only original is at the
left boundary, its maximum support is $S_L=\{0,\ldots,r\}$. Let
$U=\{r+1,\ldots,n-1\}$, so $|U|=n-r-1$. Choose a base score vector with
$s_0=1$ and all other scores distinct and below one. For example,
$s_i=-i$ for $i\ge1$ is exactly representable when $n\le2^{53}$.
At every width $w\le L$, the first admitted position is the original:
$A=1$, $B=0$, $H=K_w=1$, and the prediction count is one. Recall jumps
from zero to one at precision one and remains one thereafter. Hence every
width has PR area one, and the scalar average is one.

Suppose the algorithm's execution on this base input terminates without
inspecting some $j\in U$. Change only $s_j$ to two. The modified input is
still finite and pairwise distinct. Since $j\notin S_L$, it is outside the
support and every expanded component at every $w\le L$. On the modified
input it is admitted first, with zero recall; the original is admitted
second, with $A=1$, $B=0$, $H=K_w=1$, and prediction count two. The only
nonzero recall increment is therefore from zero to one at precision
$1/2$. All later recall increments vanish. Every width has PR area $1/2$,
and so does their scalar average.

The two adaptive query transcripts are identical. Inductively, identical
previous query answers make a deterministic algorithm choose the same next
query or termination action; every queried coordinate has the same value
because coordinate $j$ was never read in the base execution. Consequently
the algorithm returns the same output on both inputs, contradicting their
different exact scalar PR values. It must therefore inspect every coordinate
in $U$ on the base input. Counting these coordinates gives the stated bound.
An algorithm returning both scalar integrals must return the PR component,
so the same adversary applies to that output pair.
\end{proof}

\begin{corollary}[Linear-time optimal singleton subclass]
\label{cor:singleton-linear-optimal}
Under the fixed-key word-RAM and unit-cost exact metric-arithmetic model,
for inputs with $P_0=K=1$, pairwise distinct validated binary64 scores,
and $0\le L\le n$, the active-prefix evaluator uses $O(n)$ time and
$O(n)$ auxiliary words, uniformly over anomaly locations and score
orderings. Its worst-case time order matches the deterministic
$\Omega(n)$ score-inspection lower bound for exact scalar VUS-PR and for
the pair of scalar VUS-ROC and VUS-PR.
\end{corollary}
\begin{proof}
The singleton mass bound gives $J_w\le3$ for every width. Preprocessing and
all width evaluations therefore take $O(n+L)=O(n)$ time; $Q\le n$ gives
the memory bound. The boundary-location inputs in
Theorem~\ref{thm:singleton-query-lower} belong to this subclass and require
at least $n/4$ score inspections. A constant number of such inspections per
word operation gives the matching worst-case time lower bound.
\end{proof}

The lower bound is a finite-instance query statement for feasible distinct
binary64 inputs; the explicit integer witness covers $4\le n\le2^{53}$.
The asymptotic notation uses the same fixed-key operation-count convention
as the upper bound, rather than asserting an unbounded sequence of distinct
values in a literally finite 64-bit universe. Input validity need not be
verified for the lower bound. No assumption promises that the original
has the highest score: that property holds only for the adversary's base
input. The conclusion is deterministic input-inspection optimality within
the stated singleton subclass. It gives no $\Omega(L)$ lower bound for
$L\gg n$, no general $\Omega(n+L)$ lower bound, no randomized lower bound,
and no firstness or practical-performance claim.
